\documentclass[11pt]{article}
\usepackage[T1]{fontenc}
\usepackage{lmodern,amsmath,amssymb,amsthm,booktabs}
\usepackage[margin=1in]{geometry}
\usepackage[colorlinks=true,allcolors=blue]{hyperref}
\newtheorem{theorem}{Theorem}
\newcommand{\F}{\mathbb F}
\title{A four-vertex counterexample to Knudson's\protect\\ persistence-gradient conjecture}
\author{Alec Kriebel\thanks{Independent researcher. ORCID: \href{https://orcid.org/0009-0001-9320-500X}{0009-0001-9320-500X}.}}
\date{23 September 2026 \quad (version 1.0.0)}
\begin{document}
\maketitle
\begin{abstract}
Knudson conjectured that every nonincident persistence pair is joined by
exactly one gradient path in the discrete vector field obtained by retaining
the incident persistence pairs. We give a simplex-wise filtration of a
triangulated interval with four vertices and three edges for which one such
pair has no gradient path. This disproves Conjecture 2 in his contribution
to Oberwolfach Report 29/2008. The proof consists of three boundary reductions
and an inspection of the matching. The same example works over any field.
\end{abstract}

\section{The statement and the example}
Let $P$ be the persistence pairing of a finite simplicial complex whose
filtration adds one simplex at a time, with coefficients in $\F_2$.
For a codimension-one face relation $\sigma<\tau$, put
\[
 V_P=\{(\sigma,\tau)\in P:\sigma<\tau\}.
\]
Knudson's Conjecture 2 \cite[p.~1629]{Knudson} predicts a unique gradient
path in this fixed field for every nonincident pair in $P$.
We use the directed incidence graph: unmatched incidences point from a
simplex to a facet, and matched incidences point upward. A gradient path
from a critical edge starts at an endpoint and then follows matched upward
steps alternating with downward steps to the other endpoint.

Take vertices $a,b,c,d$ and edges $cd,bd,ac$, with no higher simplices.
Here $xy$ denotes $\{x,y\}$. The underlying space is the interval
$a\! -\! c\! -\! d\! -\! b$.
Insert the seven nonempty simplices in the order
\begin{equation}\label{eq:order}
             a,\ b,\ c,\ d,\ cd,\ bd,\ ac.
\end{equation}
Every face appears before its coface, so this is an admissible filtration.
The indices $1,\ldots,7$ are filtration values, not a discrete Morse function
whose own gradient is claimed to equal $V_P$.

\begin{theorem}
For the filtration \eqref{eq:order}, the finite persistence pairs and their
incident submatching are
\[
 P=\{(d,cd),(c,bd),(b,ac)\},\qquad V_P=\{(d,cd)\}.
\]
The pair $(b,ac)$ is nonincident, and there are zero gradient paths from
$ac$ to $b$ in $V_P$. Consequently Knudson's Conjecture 2 is false.
\end{theorem}

\begin{proof}
Order the vertex rows as $a<b<c<d$ and edge columns as $cd<bd<ac$.
Over $\F_2$, the usual left-to-right persistence reduction \cite{ELZ}
eliminates repeated latest nonzero rows using earlier reduced columns:
\begin{align*}
R(cd)&=\partial(cd)=c+d,\\
R(bd)&=\partial(bd)+R(cd)=(b+d)+(c+d)=b+c,\\
R(ac)&=\partial(ac)+R(bd)=(a+c)+(b+c)=a+b.
\end{align*}
The three pivots are respectively $d,c,b$, giving precisely the stated
pairs. The vertex $a$ is unpaired. Of those pairs, only $d<cd$ is an
incidence, giving the stated $V_P$.

For an independent check of the pairing, $cd$ kills the component born at
$d$; $bd$ merges the component born at $c$ into the older one born at $b$;
and $ac$ merges that surviving component into the oldest one, born at $a$.
This is the same pairing in ordinary $H_0$.

The complete directed incidence graph has the arrows
\[
 ac\to a,\quad ac\to c,\quad bd\to b,\quad bd\to d,
 \quad d\to cd,\quad cd\to c.
\]
It is acyclic. In particular $V_P$ is a discrete gradient, and both $b$
and $ac$ are critical. Either first step from $ac$ reaches $a$ or $c$;
both are unmatched vertices, so the path stops. The reachable set is
exactly $\{ac,a,c\}$, which excludes $b$.
\end{proof}

\section{Why cancellation does not repair the fixed field}
The reduction of $ac$ uses $R(bd)$ to remove its pivot at $c$. But the
persistence pair $(c,bd)$ is nonincident and contributes no upward step
to $V_P$. Algebraic column elimination therefore need not trace a path
in that field.

There is a unique path for the earlier pair $(c,bd)$:
\[
                  bd\to d\to cd\to c.
\]
Reversing it gives the new matching
$V'=\{(d,bd),(c,cd)\}$, in which the path
\[
                  ac\to c\to cd\to d\to bd\to b
\]
now exists. This observation concerns $V'$, whereas the conjecture specifies
$V_P$. It is consistent with the role of ordered cancellations in the
surface simplification results of Bauer, Lange, and Wardetzky
\cite[Lemma~9]{BLW}; we claim no counterexample to those results.

\paragraph{Other coefficients.}
Orient each edge from its alphabetically earlier to its later vertex.
Over any field the identical elimination, with subtraction, gives
\[
 R(cd)=d-c,\qquad
 R(bd)=(d-b)-(d-c)=c-b,\qquad
 R(ac)=(c-a)-(c-b)=b-a.
\]
All three pivot coefficients are $1$. The pairs, matching, and absence of
a path are unchanged. No minimality claim is needed or made.

\paragraph{Reproducibility and status.}
The companion package provides two standard-library Python verifiers:
one uses boundary reduction and union-find; the other reconstructs the
$H_0$ rank invariant from component partitions and also checks signed
reduction and path reversal. The proof above is independent of these
programs. A documented literature audit completed on 23 September 2026
located no earlier refutation of the stated conjecture. This is a bounded
search, not a guarantee of priority; in particular, an author bibliography
lists related work in preparation. The audit records inaccessible sources.
This unrefereed note and its checks were prepared with OpenAI Codex and
separate AI audit agents, from a candidate construction supplied by the
author from prior AI-assisted work.

\begin{thebibliography}{9}
\bibitem{Knudson}
K.~P. Knudson, Persistent homology and discrete Morse theory,
in \emph{Computational Algebraic Topology}, Oberwolfach Report 29/2008,
\emph{Oberwolfach Reports} \textbf{5} (2008), 1628--1630.
Conjecture 2, p.~1629.
\href{https://doi.org/10.4171/OWR/2008/29}{doi:10.4171/OWR/2008/29}.

\bibitem{ELZ}
H. Edelsbrunner, D. Letscher, and A. Zomorodian,
Topological persistence and simplification,
\emph{Discrete Comput. Geom.} \textbf{28} (2002), 511--533.
\href{https://doi.org/10.1007/s00454-002-2885-2}{doi:10.1007/s00454-002-2885-2}.

\bibitem{BLW}
U. Bauer, C. Lange, and M. Wardetzky,
Optimal topological simplification of discrete functions on surfaces,
\emph{Discrete Comput. Geom.} \textbf{47} (2012), 347--377.
\href{https://doi.org/10.1007/s00454-011-9350-z}{doi:10.1007/s00454-011-9350-z}.
\end{thebibliography}
\end{document}
